\begin{abstract}
  Энэ дипломын ажил нь Монгол Улсын Их Сургуулийн бакалаврын түвшиний speaker recognition ба speaker diarization-т
  хүлээн зөвшөөрөгдсөн стандартын дагуу боловсруулсан судалгааны ажил юм. Энэхүү бүтээл нь дараах асуудлуудыг харуулна:
  - хэл таних системийн архитектурын оновчтой шийдлүүдийг боловсруулж турших;
  - олон speakers-тай нөхцөлд яриаг тусгаарлах diarization-ийн үр ашигтай аргачлал боловсруулах;
  - x-vector/PLDA зэрэг орчин үеийн үгсийн төлөв байдлын үнэлгээ, өгөгдөл бэлтгэл, үнэлгээ, үр дүнгийн шинжилгээ.
  
  Энэ ажил нь өгөгдөл боловсруулалт, эмбэддокын гаргац, үнэлгээ, үр дүнгийн анализыг хамарсан бүтэн туршлагатой судалгааны тайлан юм. 
  Эх сурвалжийн өгөгдлийг туршилтад ашиглах замаар DER/ EER зэрэг метрикүүдээр үнэлнэ.
  
  Тайлан нь дараах бүтэцтэй: онол, өгөгдөл ба аргачлал, туршилт, үр дүн, хэлэлцүүлэг ба ирээдүйн ажил.
\end{abstract}
